\documentclass[a4paper, twoside]{report}
\usepackage[T1]{fontenc} % Font encoding
\usepackage[utf8]{inputenc} % Input encoding
\usepackage[english,french]{babel} % Main language French, some English bits and pieces
\usepackage[backend=biber, style=verbose-ibid, sorting=nyt, citepages=omit]{biblatex} % For Bibliography
\usepackage{graphicx} % For external graphics etc
\usepackage{lipsum} % To generate Lorem Ipsums
\usepackage[a4paper,top=3cm,bottom=3cm,left=3cm,right=3cm]{geometry} % Margins and paging
\usepackage[fontsize=13pt]{scrextend} % Bigger Font
\usepackage{hyperref} % Clickable ToC and hyperlinks
\usepackage{listings} % Insert code snippet with clear formatting
\usepackage{appendix} % Better Appendix formatting
\raggedbottom 
\pagestyle{headings}

%\renewcommand{\appendixtocname}{Annexes} % Making sure the appendix' names are correctly translated
%\renewcommand{\appendixpagename}{\appendixtocname}

\usepackage{xcolor} % Use for custom colors, import it last because of collisions with other packages (eg. TikZ)

\hypersetup{
    colorlinks,
    linkcolor={red!50!black},
    citecolor={blue!50!black},
    urlcolor={blue!80!black}
}

\begin{document}
\pagenumbering{roman}


\begin{titlepage}
\begin{figure}[!htb]
    \centering
    \includegraphics[keepaspectratio=true,scale=0.25]{./au.eps}
\end{figure}

\begin{center}
    \LARGE{AVIGNON UNIVERSITÉ}
    \vspace{5mm}
    \\ \large{Département STS - CERI - Licence Informatique}
    \vspace{5mm}
    \\ \LARGE{UE Administration des systèmes d'exploitation}
\end{center}

\vspace{15mm}
\begin{center}
    {\LARGE{\bf Mini projet\\\vspace{5mm} Administration des Systèmes d’Exploitation}}
\end{center}
\vspace{30mm}

\begin{minipage}[t]{0.47\textwidth}
	{\large{Enseignant:}{\normalsize\vspace{3mm}
	\bf\\ \large{JABAIAN Bassam \vspace{2mm}\\ NOCERA Pascal }}}
\end{minipage}
\hfill
\begin{minipage}[t]{0.47\textwidth}\raggedleft
	{\large{Auteur:}{\normalsize\vspace{3mm}
    \bf\\ \large{SALLOUM Guillaume}}}
\end{minipage}

\vspace{20mm}
\hrulefill
\\\centering{\large{Année 2022/2023, Semestre 4}}
\\\vspace{2em}\centering{\small{Template d'Alexandre Klein, distribué sous Licence GNU GPLv3}}

\end{titlepage}

\clearpage


{ \hypersetup{hidelinks} \tableofcontents \addcontentsline{toc}{section}{\numberline{}Table des matières} }

\setcounter{page}{1}\pagenumbering{arabic}

\lstdefinestyle{bash}{
    commentstyle=\color{green!40!black},
    keywordstyle=\color{black!70!black},
    basicstyle=\footnotesize\ttfamily,
    columns=flexible,
    language=bash,
    belowcaptionskip=1\baselineskip,
    extendedchars=false,
    breaklines=true,
    numbers=left,
    numberstyle=\tiny\color{gray},
    frame=lines,
    inputencoding=utf8,
}

\lstdefinestyle{python}{
    commentstyle=\color{green!40!black},
    keywordstyle=\color{black!70!black},
    basicstyle=\footnotesize\ttfamily,
    breakatwhitespace=false,         
    breaklines=true,                 
    captionpos=b,                    
    keepspaces=true,                 
    numbers=left,
    numberstyle=\tiny\color{gray},                 
    showspaces=false,                
    showstringspaces=false,
    showtabs=false,                  
    tabsize=2,
    frame=lines,
    language=python,
    inputencoding=utf8,
}

\chapter{Collecte d'informations}

\begin{lstlisting}[style=python]
# Sonde CPU
#!/bin/python

import psutil

print("Informations CPU :")
print(psutil.cpu_percent(0.1))
print("Fin affichage CPU")
\end{lstlisting}

%\includegraphics[keepaspectratio=true,scale=0.60]{./graphics/images/sonde_cpu.png}


\begin{lstlisting}[style=bash]
#!/bin/bash

# Lister les processus de l'utilisateur uapv2103100
# Sonde processus
echo "Sonde processus"

ps -eF | grep uapv210+

# Fin affichage processus

echo "Fin affichage processus"

\end{lstlisting}

%\includegraphics[keepaspectratio=true,scale=0.60]{./graphics/images/processus.png}

\begin{lstlisting}[style=bash]
#!/bin/bash

# Afficher les espaces du disques utilises et les espaces libres

echo "Sonde disque"
du -shk /home/*;
echo "Fin affichage disque"

\end{lstlisting}

%\includegraphics[keepaspectratio=true,scale=0.60]{./graphics/images/sonde_disk.png}

\newpage
\begin{lstlisting}[style=bash]
#!/bin/bash

# Sonde globale qui appelle toutes les sondes

echo "Sonde CPU"
python3 sonde1.py

echo "Sonde processus"
bash sonde2.sh

echo "Sonde disque"
bash sonde3.sh

\end{lstlisting}

\begin{lstlisting}[style=python]
import psutil

print("Memoire par utilisateur")
users = []
mem = []
for proc in psutil.process_iter():
try:
    pinfo = proc.as_dict(attrs=['username', 'memory_percent'])
    users.append(pinfo['username'])
except (psutil.NoSuchProcess, psutil.AccessDenied, psutil.ZombieProcess):
    pass

users2 = []
for user in users:
if user not in users2:
    users2.append(user)
#print(users)
#print(users2)
for j in users2:
mem.append(0)
k = 0
for j in users2:
for proc in psutil.process_iter():
    try:
        pinfo = proc.as_dict(attrs=['username', 'memory_percent'])
        if (pinfo['username'] == j):
            mem[k] = mem[k] + pinfo['memory_percent']
    except (psutil.NoSuchProcess, psutil.AccessDenied, psutil.ZombieProcess):
        pass
print(j, "utilise", mem[k], "% de la mémoire")
k = k + 1
print("Fin")

\end{lstlisting}

\chapter{Stockage et archivage}

\begin{lstlisting}[style=bash]
#!/bin/bash

# Stockage
i=1

# Effacer les donnees precedentes pour actualiser les donnees
# rm ./donnees.txt

while true
do
    echo "Affichage des sondes";
    (python3 sonde1.py && ./sonde2.sh && ./sonde3.sh && python3 sonde4.py) | tee -a ./donnees.txt;
    ((i++));
    sleep 5;
done

\end{lstlisting}

\begin{lstlisting}[style=python]
#!/bin/python
# Formatter l'affichage des sondes

liste=[]
disk=[]
cpu=[]
proc=[]
users=[]
mem=[]
mem2=[]

with open("donnees.txt", "r") as f:
    liste=[line.strip() for line in f]

for i in range(0, len(liste)):
        if (liste[i] == "Informations CPU :"):
                while (liste[i+1] != "Fin affichage CPU"):
                        head, sep, tail = liste[i+1].partition(':')
                        cpu.append(int(float(head)))
                        i = i + 1

for i in range(0, len(liste)):
    if (liste[i] == "Sonde disque"):
        while (liste[i+1] != "Fin affichage disque"):
            head, sep, tail = liste[i+1].partition('\t/home/')
            disk.append(head)
            i = i + 1

for i in range(0, len(liste)):
    if (liste[i] == "Sonde processus"):
        while (liste[i+1] != "Fin affichage processus"):
            head, sep, tail = liste[i+1].partition('\t/home/')
            proc.append(head)
            i = i + 1

for i in range(0, len(liste)):
    if (liste[i] == "Memoire par utilisateur"):
        while (liste[i+1] != 'Fin'):
            head, sep, tail = liste[i+1].partition(' utilise ')
            users.append(head)
            mem.append(tail)
            i = i + 1

\end{lstlisting}

\newpage
\begin{lstlisting}[style=python]
import mysql.connector
from mysql.connector import Error
from format import liste, disk, cpu, proc
from sonde4 import users2 as users, mem
import xml.etree.ElementTree as ET
import xml.dom.minidom
import time

try:
    connection = mysql.connector.connect(
        host='localhost',
        database='sondes',
        user='uapv2103100',
        password='1234')

    if connection.is_connected():
        curseur = connection.cursor()
        bdd_info = connection.get_server_info()
        i = 0
        TEMPS = int(time.time())
        k = TEMPS
        while (i < len(users)):
            info ="""INSERT INTO sonde
                        (Temps, CPU, Processus, Disque, Memoire)
                        VALUES (%s,%s,%s,%s,%s)"""
            values = (TEMPS,cpu[i],proc[i],disk[i],mem[i])
            curseur.execute(info, values)
            connection.commit()
            print(curseur.rowcount, "Insertion de valeurs dans la table sonde")
            i = i + 1
            TEMPS = TEMPS + 1
        curseur.close()

        # Gestionnaire d'historique et backup en XML
        curseur = connection.cursor()
        j = 0
        root = ET.Element('root')
        doc = ET.SubElement(root, 'doc')
        while (j < len(users)):
            curseur.execute('SELECT * FROM sonde')
            for row in curseur:
                print(row)
            u = ET.SubElement(doc, 'user')
            ET.SubElement(u, "TIME", name="TIME").text = "%s" % row[0]
            ET.SubElement(u, "CPU", name="CPU").text = "%s" % row[1]
            ET.SubElement(u, "PROC", name="PROC").text = "%s" % row[2]
            ET.SubElement(u, "DISK", name="DISK").text = "%s" % row[3]
            ET.SubElement(u, "RAM", name="RAM").text = "%s" % row[4]
            j = j + 1
        tree = ET.ElementTree(root)
        tree.write("db.xml")
        dom = xml.dom.minidom.parse("db.xml")
        p = dom.toprettyxml(indent = '\t')
        f = open("db.xml", "r+")
        f.truncate(0)
        f.write(p)
        f.close()
        curseur.close()

        curseur = connection.cursor()
        cmd = "DELETE FROM sonde ORDER BY Temps LIMIT 10"
        curseur.execute(cmd)
        print("Base nettoyee")
        curseur.close()

except Error as e:
    print("Erreur", e)
finally:
    if connection.is_connected():
        curseur.close()
        connection.close()
        print("Fermeture de la connextion a MySQL")

\end{lstlisting}


\chapter{Affichage et Alerte}

\begin{lstlisting}[style=python]
from db import users, disk, mem, cpu
import pygal
from pygal.style import DarkStyle
from pygal.style import NeonStyle
import smtplib
import getpass
import time

f = open("donnees.txt", "r")
data = f.read()
print(data)

# Condition alerte
utilisation_ram = sum(mem)
alerte = 0
if (utilisation_ram > 25):
print ("RAM : ", utilisation_ram, "%")
print("Alerte RAM")
alerte = 1

# Envoyer un mail
if (alerte == 1):
password = getpass.getpass(prompt = 'Mot de passe : ')
smtp = smtplib.SMTP_SSL('smtpz.univ-avignon.fr:465')
smtp.login("guillaume.salloum@alumni.univ-avignon.fr", password)
msg = """Subject: Alerte RAM 
L utilisation de la RAM sur la VM est superieure a 25% ! """
smtp.sendmail("guillaume.salloum@alumni.univ-avignon.fr", "guillaume.salloum@alumni.univ-avignon.fr", msg)

# Graphiques
chart = pygal.HorizontalLine(title = "Evolution du %CPU", style=NeonStyle)
chart.x_labels = users
chart.add('CPU', int(cpu[0]))
chart.render_to_file('SondeCPU.svg')

chart2 = pygal.StackedLine(title = "Evolution de l espace utilise du disque", style=DarkStyle)
chart2.x_labels = users
chart2.add('Disque', int(disk[1]))
chart2.render_to_file('SondeDISK.svg')

chart3 = pygal.Bar(title = 'Evolution de la RAM/Utilisateur', style=NeonStyle)
chart3.x_labels = users
chart3.add('RAM', mem)
chart3.render_to_file('SondeRAM.svg')

\end{lstlisting}

\chapter{Interface Web}

\begin{lstlisting}[style=python]
#!/bin/bash

i=1;
#rm donnees.txt
while true
do
echo "Affichage des sondes";
(python3 sonde1.py && ./sonde2.sh && ./sonde3.sh && python3 sonde4.py) | tee -a ./donnees.txt;
((i++));
python3 format.py
python3 db.py
python3 alerte.py
sudo cp SondeCPU.svg ./Sonde1.svg
sudo cp SondeDISK.svg ./Sonde2.svg
sudo cp SondeRAM.svg ./Sonde3.svg
sleep 10;
done

\end{lstlisting}

\lstdefinelanguage{HTML5}{
    sensitive=true,
    keywords={%
    % JavaScript
    typeof, new, true, false, catch, function, return, null, catch, switch, var, if, in, while, do, else, case, break,
    % HTML
    html, title, meta, style, head, body, script, canvas,
    % CSS
    border:, transform:, -moz-transform:, transition-duration:, transition-property:,
    transition-timing-function:
    },
    % http://texblog.org/tag/otherkeywords/
    otherkeywords={<, >, \/},   
    ndkeywords={class, export, boolean, throw, implements, import, this},   
    comment=[l]{//},
    % morecomment=[s][keywordstyle]{<}{>},  
    morecomment=[s]{/*}{*/},
    morecomment=[s]{<!}{>},
    morestring=[b]',
    morestring=[b]",    
    alsoletter={-},
    alsodigit={:}
}

\lstdefinestyle{html}{
    % Basic design
    backgroundcolor=\color{lightgray},
    basicstyle={\small\ttfamily},   
    frame=l,
    % Line numbers
    xleftmargin={0.75cm},
    numbers=left,
    stepnumber=1,
    firstnumber=1,
    numberfirstline=true,
    % Code design
    identifierstyle=\color{black},
    keywordstyle=\color{blue}\bfseries,
    ndkeywordstyle=\color{greenCode}\bfseries,
    stringstyle=\color{ocherCode}\ttfamily,
    commentstyle=\color{darkgray}\ttfamily,
    % Code
    language={HTML5},
    tabsize=2,
    showtabs=false,
    showspaces=false,
    showstringspaces=false,
    extendedchars=true,
    breaklines=true
}

\begin{lstlisting}[style=html]
<!DOCTYPE html>
<html>
<head>
<title>Projet UE Administration des Systèmes d'Exploitation 2022-2023</title>
</head>
<body>

<h1 style="text-align:center">Projet UE Administration des Systèmes d'Exploitation 2022-2023</h1>
<div style="text-align:center"></div>
<button type="button" style="text-align:center" onclick="toggle()">Sonde CPU</button>
<button type="button" style="text-align:center" onclick="toggle2()">Sonde Disque</button>
<button type="button" style="text-align:center" onclick="toggle3()">Sonde RAM</button>
</div>
<div style="text-align:center" id="1">
<embed type="image/svg+xml" src="Sonde1.svg" width="1000" height="600" />
</div>
<div style="text-align:center" id="2">
<embed type="image/svg+xml" src="Sonde2.svg" width="1000" height="600" />
</div>
<div style="text-align:center" id="3">
<embed type="image/svg+xml" src="Sonde3.svg" width="1000" height="600" />
</div>
<script>
function toggle() {
    var x = document.getElementById("1");
    if (x.style.display === "none") {
    x.style.display = "block";
    } else {
    x.style.display = "none";
    }
}

function toggle2() {
    var x = document.getElementById("2");
    if (x.style.display === "none") {
    x.style.display = "block";
    } else {
    x.style.display = "none";
    }
}

function toggle3() {
    var x = document.getElementById("3");
    if (x.style.display === "none") {
    x.style.display = "block";
    } else {
    x.style.display = "none";
    }
}
</script>

</body>
</html>
\end{lstlisting}



\end{document}
